\documentclass[../../../include/open-logic-section]{subfiles}

\begin{document}

\olfileid{sth}{choice}{tarskiscott}
\olsection{టార్స్కీ--స్కాట్ యుక్తి}

\olref[cardinals][cardsasords]{defcardinalasordinal}లో కార్డినల్ సంఖ్యలను క్రమసంఖ్యలుగా నిర్వచించాం. అందుకోసం సుక్రమపరచడం స్వయంసిద్ధాన్ని అనుకున్నాం. అలా చేసిన ఏకైక కారణం అది ``సాధారణంగా అనుసరించే పద్ధతి'' కావడమే.

కాబట్టి సుక్రమపరచడం చుట్టూ ఉన్న తాత్త్విక ప్రశ్నలను చర్చించే ముందు, ఆ సాధారణ పద్ధతికి భిన్నంగా వెళ్లి, సుక్రమపరచడాన్ని అనుకో\emph{కుండానే} కార్డినల్ సంఖ్యల సిద్ధాంతాన్ని అభివృద్ధి చేయ\emph{వచ్చని} స్పష్టం చేసుకోవాలి. \olref[sfr][siz][]{chap}లో ఇచ్చిన $\cardeq{A}{B}$, $\cardle{A}{B}$, $\cardless{A}{B}$ నిర్వచనాలను అలాగే వాడవచ్చు. \emph{కార్డినల్ సంఖ్య}కు మాత్రం కొత్త భావన కావాలి.

సరళంగా ఆలోచిస్తే $A$ కార్డినాలిటీని ఇలా నిర్వచించడానికి ప్రయత్నించవచ్చు:
\[
	\Setabs{x}{\cardeq{A}{x}}.
\]
దీన్ని \olref[cardinals][hp]{sec} చివర్లో సూచించిన $\# x Fx$కు ఫ్రేగే ఇచ్చిన నిర్వచనంతో పోల్చవచ్చు. అక్కడ సూచించిన కారణాలవల్లే ఈ నిర్వచనం విఫలమవుతుంది. ప్రతి ఏకమూలక సమితీ $\{\emptyset\}$తో సమసంఖ్యకం. కానీ స్థాయిక్రమంలోని ప్రతి ఉత్తరవర్తి దశలో కొత్త ఏకమూలక సమితులు ఏర్పడతాయి (ముందరి దశనే ఏకమూలక సమితిగా తీసుకోండి). అందువల్ల $\Setabs{x}{\cardeq{A}{x}}$కు స్థాయి ఉండదు; అది సమితిగా లేదు.

ఈ సమస్యను దాటడానికి టార్స్కీ, స్కాట్‌ల యుక్తిని వాడతాం:\footnote{గుర్తుచేసుకోండి: వేరుగా స్పష్టంగా చెప్పకపోతే అన్ని సూత్రాల్లో పరామితులు ఉండవచ్చు.}

\begin{defn}[టార్స్కీ--స్కాట్]
ఏ సూత్రం $\phi(x)$కైనా, $[ x : \phi(x)] $ అనేది $\phi(x)$ను తృప్తిపరిచే, సాధ్యమైన అత్యల్ప స్థాయి గల అన్ని $x$ల సమితి (అలాంటి $\phi$లు లేకపోతే $\emptyset$).
\end{defn}

ఈ నిర్వచనం సమంజసమో తనిఖీ చేద్దాం. $\ZF$లో పనిచేస్తే, ప్రతి $x$కు $\setrank{x}$ ఉందని \olref[spine][foundation]{zfentailsregularity} నిర్ధారిస్తుంది. ఇప్పుడు $\phi$ను తృప్తిపరిచే వస్తువులు ఏవైనా ఉంటే, $(\exists x\subseteq V_\alpha)\phi(x)$ అయ్యే అత్యల్ప స్థాయి $\alpha$ను తీసుకోవచ్చు; అంటే $(\forall \beta
\in \alpha)(\forall x \subseteq V_\beta)\lnot \phi(x)$. అప్పుడు విభజన ద్వారా $[x : \phi(x)]$ను $\Setabs{x \in
V_{\alpha+1}}{\phi(x)}$గా నిర్వచించవచ్చు.

టార్స్కీ--స్కాట్ యుక్తి సమర్థించబడింది కాబట్టి, ఇప్పుడు దానితో కార్డినాలిటీ భావనను నిర్వచించవచ్చు:

\begin{defn}
$A$ యొక్క \textsc{ts}-కార్డినాలిటీ $\text{tsc}(A) = [x :
\cardeq{A}{x}]$.
\end{defn}

\textsc{ts}-కార్డినల్ నిర్వచనం సుక్రమపరచడాన్ని వాడదు. అయినా ఆ స్వయంసిద్ధం లేకుండానే \emph{\textsc{ts}-కార్డినల్ సంఖ్యలు}, \olref[cardinals][cardsasords]{defcardinalasordinal}లో నిర్వచించిన \emph{కార్డినల్ సంఖ్యల} వలె ప్రవర్తిస్తాయని చూపవచ్చు. ఉదాహరణకు \olref[cardinals][cardsasords]{lem:CardinalsBehaveRight}, \olref[card-arithmetic][opps]{lem:SizePowerset2Exp}లను \textsc{ts}-కార్డినల్ సంఖ్యల భాషలో మళ్లీ రాస్తే, వాటి నిరూపణలు సుక్రమపరచడాన్ని అనుకోకుండానే $\ZF$లో పనిచేస్తాయి.

ఇదే సందర్భంలో టార్స్కీ--స్కాట్ యుక్తితో క్రమసంఖ్యల సిద్ధాంతాన్నీ అభివృద్ధి చేయవచ్చని గమనించాలి. $\tuple{A, <}$ సుక్రమం అయితే $\text{tso}(A, <) = [\tuple{X,
R} : \ordeq{\tuple{A, <}}{\tuple{X, R}}]$గా తీసుకుందాం. కార్డినల్ సంఖ్యలు, క్రమసంఖ్యలపై ఈ పద్ధతి గురించి మరింత తెలుసుకోవడానికి \citet[chs.~9--12]{Potter2004} చూడండి.

\end{document}
